\scanpage{78}
% archive sheet number: 77
% L1-6
Soit $G$ un groupe où opère $\mathbb Z$, \inserted{${\unsure{\Pi|T}}^1=T$.} Pour $\forall$ entier $n$, soit $G_n$ l'ens des éléments de $G$ tels que $\Pi^n x=x$.
On suppose que $G=\bigcup G_n$ (les $G_n$ forment une suite croissante). Alors $G_n$ est un $\mathbb Z_n$-groupe, et on peut considérer sa cohomologie \inserted{$H^i$} (\unsure{si $i>1$, supposons que} $G$ commutatif), alors on a
% L7-11
\[
H^0(\mathbb Z_n,G_n)=G_1\qquad\text{ens. des \unsure{él$^{ts}$|él$^t$} de $G$ invariants par $\mathbb Z$}
\]
\[
\begin{aligned}
H^1(\mathbb Z_n,G_n)&=\text{ens. des éléments $\alpha$ de $G_n$ tels que}\\
&\quad\alpha\,T\alpha\,T^2\alpha\cdots T^{n-1}\alpha=1,\quad\text{mod la relation d'équivalence}\\
&\quad\text{définie par }\beta\simeq\alpha\text{ si et seulement si}\\
&\quad\text{il existe }\unsure{u|x}\in G_n\text{ tel que }\beta=u^{-1}\alpha\,Tu
\end{aligned}
\]
% L12-14
D'ailleurs, supposons que $n\mid m$, alors $\mathbb Z_n$ est un quotient de $\mathbb Z_m$, \deleted{\unsure{et on a}} par un sous-groupe $N_{n,m}$, et on a
\[
\boxed{G_n=G_m^{N_{nm}}}.
\]
On a donc fonctoriellement une suite exacte
% L15
\diagram{
\begin{tikzcd}[ampersand replacement=\&]
e\arrow[r]\& H^1(\mathbb Z_n,G_n)\arrow[r,"\text{inj}"]\&H^1(\mathbb Z_m,G_m)\arrow[r]\&H^1(\underbrace{N_{nm}}_{\substack{\Vert\\\mathbb Z_{m/n}}},G_m)^{\unsure{\mathbb Z_m|\mathbb Z_n}}
\end{tikzcd}}
% L16-20
[Si \deleted{$\unsure{G_n}$} $G$ est abélien, on a aussi $H^i(\mathbb Z_n,G_n)$, et on a des hom
\[
H^i(\mathbb Z_n,G_n)\longrightarrow H^i(\mathbb Z_m,G_m)
\]
d'ailleurs, on a à partir de $i=1$, une \uline{périodicité} de période 2 dans les $H^i(\mathbb Z_n,G_n)$, \deleted{\unsure{pour les}} \unsure{les $i>1$ n'y gagnent pas grand chose}].

% L21-23
\uline{Lemme!} Soit $\alpha\in G$ tel que $\alpha\,T\alpha\,T^2\alpha\cdots T^{n-1}\alpha=1$, alors $\alpha\in G_n$.
Supposons $\alpha,\beta\in G$ satisfaisant aux conditions précédentes, et $\beta=u^{-1}\alpha\,Tu$ $(u\in G)$, alors $u\in G_n$.
% L24-27
En effet, \unsure{on aura}, appliquant $T$ \deleted{$\unsure{\cdot}$}
\[
T\alpha\,T^2\alpha\cdots T^{n-1}\alpha\,T^n\alpha=1
\]
i.e. $T^n\alpha$ est un \uline{inverse} \deleted{\illegible{1}} de $T\alpha\,T^2\alpha\cdots T^{n-1}\alpha$ ; comme il en est de $=$ de $\alpha$, on a $T^n\alpha=\alpha$.
De $\unsure{\equiv|=}$, écrivons $\beta\,T\beta\cdots T^{n-1}\beta=1$, en y remplaçant $\beta$ par $u^{-1}\alpha\,Tu$, on trouve
